% !PRINTED 197
\weberpage{197}
%<<CONT
is. It is easy to see that this case, if $\Theta$ has a finite
degree, cannot occur.

First suppose that $\beta$ is zero or $\gamma = 0$. It follows from (7) that
$\alpha = \delta$ must hold, and both can be assumed $= 1$.

Now if $\Theta$ is one of the two substitutions
\[
\begin{pmatrix} 1, & 0 \\ \gamma, & 1 \end{pmatrix}, \quad
\begin{pmatrix} 1, & \beta \\ 0, & 1 \end{pmatrix}
\]
then for every $\lambda$:
\[
\Theta^\lambda = \begin{pmatrix} 1, & 0 \\ \lambda\gamma, & 1 \end{pmatrix} \text{ or } \begin{pmatrix} 1, & \lambda\beta \\ 0, & 1 \end{pmatrix},
\]
as one easily finds by complete induction. Therefore,
since $\beta$ and $\gamma$ cannot both be zero if $\Theta$ is not
the identical substitution, $\Theta^\lambda$ cannot equal 1 for any positive $\lambda$,
as it would have to if $\lambda$ were equal to the degree of
$\Theta$.

But if $\beta$ is different from zero, then, assuming now
$\Delta = 1$ and, by (8), $\alpha + \delta = 2$, we set
\[
L = \begin{pmatrix} \beta, & 0 \\ 1-\alpha, & \beta^{-1} \end{pmatrix},
\]
and obtain
\[
\begin{aligned}
L^{-1}\Theta L &= \begin{pmatrix} \beta^{-1}, & 0 \\ \alpha-1, & \beta \end{pmatrix} \begin{pmatrix} \alpha, & \beta \\ \gamma, & \delta \end{pmatrix} \begin{pmatrix} \beta, & 0 \\ 1-\alpha, & \beta^{-1} \end{pmatrix} \\
&= \begin{pmatrix} 1, & \beta^{-1} \\ 0, & 1 \end{pmatrix},
\end{aligned}
\]
hence
\[
\Theta = L \begin{pmatrix} 1, & \beta^{-1} \\ 0, & 1 \end{pmatrix} L^{-1},
\]
and from this
\[
\Theta^\lambda = L \begin{pmatrix} 1, & \lambda\beta^{-1} \\ 0, & 1 \end{pmatrix} L^{-1},
\]
which again cannot equal 1 for any positive $\lambda$. We
therefore have the theorem:

\par\noindent 1.\ \gesperrt{Each of the $n-1$ non-identical substitutions
of a group of degree $n$ has two poles different
from one another.}

Every non-identical substitution of the group $G$ thus gives us
two poles. But one and the same value can occur as a pole
in several different substitutions. If we count
one of these values $h$ times when it occurs as a pole in $h$ substitutions of the
group, then the number of poles results
%CONT>>
